\section*{Chapter \ref{ch:g5:networks-first-look} --- Networks and the Internet: A First Look}

\begin{solution}{exo:g5:networks-first-look:1}
For example: a laptop (by cable or by radio), a phone (by radio, wi-fi), a
printer (by cable or by radio).
\end{solution}

\begin{solution}{exo:g5:networks-first-look:2}
No. The \omterm{def:g5:networks-first-look:internet}{Internet} is millions of \omterm{def:g5:networks-first-look:network}{networks} joined together, each with its own
\omterm{def:g1:what-is-a-computer:computer}{computers}; no single \omterm{def:g1:what-is-a-computer:computer}{computer} is the \omterm{def:g5:networks-first-look:internet}{Internet}.
\end{solution}

\begin{solution}{exo:g5:networks-first-look:3}
It asks a \omterm{def:g5:networks-first-look:server}{server} for a \omterm{def:g5:networks-first-look:page}{web page}, using the page's address, and shows the page
it receives; when we \omterm{def:g2:keyboard-and-mouse:click}{click} a \omterm{def:g5:networks-first-look:web}{link}, it asks for the next page.
\end{solution}

\begin{solution}{exo:g5:networks-first-look:4}
\omterm{def:g3:files-and-folders:file}{Files} and \omterm{def:g4:programs-and-operating-system:program}{programs}, such as \omterm{def:g5:networks-first-look:page}{web pages}, photos and messages. It sends them to
the \omterm{def:g1:what-is-a-computer:computer}{computers} that ask for them.
\end{solution}

\begin{solution}{exo:g5:networks-first-look:5}
A, then B, then D (two \omterm{def:g2:keyboard-and-mouse:click}{clicks}). From D, the only \omterm{def:g5:networks-first-look:web}{link} leads to B, and from B
the only \omterm{def:g5:networks-first-look:web}{link} leads to D: there is no way back to A.
\end{solution}

\begin{solution}{exo:g5:networks-first-look:6}
Weak: ``1234'', ``anna'', ``qwerty'' (a row of the keyboard) and ``Tom2015''
(a name and a year). Strong: ``purple-tractor-moonlight'' and ``giraffe kettle
snowstorm'' (long, three random words).
\end{solution}

\begin{solution}{exo:g5:networks-first-look:7}
$10 \times 10 \times 10 = \num{1000}$ codes, from 000 to 999.
\end{solution}

\begin{solution}{exo:g5:networks-first-look:8}
$26 \times 26 \times 26 \times 26$: $26 \times 26 = 676$, $676 \times 26 =
\num{17576}$, and $\num{17576} \times 26 = \num{456976}$ \omterm{def:g5:networks-first-look:password}{passwords}.
\end{solution}

\begin{solution}{exo:g5:networks-first-look:9}
If one website loses its \omterm{def:g5:networks-first-look:password}{passwords}, or someone sees one of them, that person
could open all the accounts that use the same \omterm{def:g5:networks-first-look:password}{password}. With a different
\omterm{def:g5:networks-first-look:password}{password} for each, only one account is in danger.
\end{solution}

\begin{solution}{exo:g5:networks-first-look:10}
By 10 for digits (each new place has 10 choices), and by 26 for small letters.
\end{solution}

\begin{solution}{exo:g5:networks-first-look:11}
Its \omterm{def:g4:programs-and-operating-system:program}{programs} visited a huge number of \omterm{def:g5:networks-first-look:page}{web pages} in advance, following their
\omterm{def:g5:networks-first-look:web}{links}, and made a list of which words appear on which page. When you search, it
only looks up your words in that list.
\end{solution}
